\section{Real Analysis Preview: Epsilon as a Contract}
\chapterline

In introductory calculus, we are told that a sequence approaches a limit:
\[
\lim_{n \to \infty} a_n = L,
\]
or that a function approaches a value:
\[
\lim_{x \to c} f(x) = L.
\]
Textbooks describe this dynamically: ``As $x$ gets closer and closer to $c$, $f(x)$ gets closer and closer to $L$.''

While that mental picture is intuitive, it lacks logical teeth. What does ``closer and closer'' actually mean? Does it mean the sequence can never drift backward? Does it mean the distance decreases at every single step? 

In the nineteenth century, mathematicians like Augustin-Louis Cauchy and Karl Weierstrass realized that if mathematics was to rest on solid bedrock, motion had to be replaced by logic. They reframed the limit not as a physical journey, but as an \textbf{ironclad adversarial contract}.

\begin{historybox}
The rigor of real analysis was born out of a scathing philosophical attack. In 1734, Bishop George Berkeley published \emph{The Analyst: A Discourse Addressed to an Infidel Mathematician}. Berkeley attacked Isaac Newton's ``fluxions'' and Gottfried Leibniz's ``differentials,'' pointing out a glaring logical contradiction in early calculus: to compute a derivative, one divided by an increment $\Delta x$ (which requires $\Delta x \neq 0$), only to declare $\Delta x = 0$ a moment later to make unwanted terms vanish. Berkeley famously scoffed:
\begin{quote}
\itshape
``And what are these same evanescent increments? They are neither finite quantities, nor quantities infinitely small, nor yet nothing. May we not call them the ghosts of departed quantities?''
\end{quote}
For nearly a century, mathematicians calculated brilliant results using these ``ghosts,'' yet lacked a defensible logical foundation. The breakthrough came when Augustin-Louis Cauchy (\emph{Cours d'Analyse}, 1821) and Karl Weierstrass in his 1860s Berlin lectures abolished infinitesimals entirely. Weierstrass realized that calculus does not require quantities that are mysteriously ``infinitely small.'' Instead, a limit is a static game played with ordinary, positive real numbers ($\eps$ and $\delta$) governed by nested quantifiers.
\end{historybox}

\subsection{The Adversarial Game: Sequences}

\begin{definition}[Limit of a Sequence]
A sequence of real numbers $(a_n)_{n=1}^\infty$ \textbf{converges} to a real number $L$, written $\lim_{n \to \infty} a_n = L$, if:
\[
(\forall \eps > 0)(\exists N \in \N)(\forall n \in \N)(n \ge N \implies |a_n - L| < \eps).
\]
\end{definition}

\subsection{Plain English}

No matter how narrow a vertical tolerance band $(L - \eps, L + \eps)$ someone draws around $L$, you can walk far enough out along the sequence—past some cutoff index $N$—that every term from $a_N$ onward stays trapped inside that band forever. Conversely, if a sequence like $a_n = (-1)^n$ keeps jumping between $+1$ and $-1$, then as soon as $\eps \le 1$, no horizontal band of half-height $\eps$ can trap both values, so terms keep escaping no matter how far out you go.

\subsection{Who Chooses What?}

To master this definition, view it as a two-player game between a \textbf{Challenger} (the skeptic) and a \textbf{Prover} (you):
\begin{enumerate}[label=\arabic*.]
\item \textbf{The Challenger's Move ($\forall \eps > 0$):} The challenger goes first and picks an error tolerance $\eps > 0$. Because they want to expose a flaw in your claim, they may choose $\eps$ as tiny as they please—say, $\eps = 0.00001$.
\item \textbf{Your Move ($\exists N \in \N$):} After seeing their $\eps$, you inspect the formula for $a_n$ and produce a finite index cutoff $N \in \N$. Because you move second, your choice of $N$ is allowed to depend on $\eps$.
\item \textbf{The Challenger's Inspection ($\forall n \ge N$):} The challenger is now free to test \emph{any} index $n \ge N$ at or past your cutoff—including the worst-case boundary term $n = N$—and you win the round only if $|a_n - L| < \eps$ holds for every single $n \ge N$:
\[
L - \eps < a_n < L + \eps.
\]
\end{enumerate}

If you have a formula or strategy that produces a winning cutoff $N(\eps)$ against \emph{every possible} positive tolerance $\eps$ the challenger can name, the sequence converges to $L$.

\begin{center}
\begin{tikzpicture}[scale=0.90, >=Latex]
  % Lane 1: Sequence Epsilon Corridor
  \begin{scope}[yshift=0cm]
    \node[anchor=west, font=\small\bfseries] at (-0.5, 3.6) {1. The Sequence Contract: $\forall \eps > 0, \exists N \in \N, \forall n \ge N, |a_n - L| < \eps$};
    \draw[rounded corners=4pt, fill=gray!3, draw=gray!25] (-0.6, -0.6) rectangle (12.2, 3.4);
    
    % Shaded tolerance band
    \fill[targetgreen!12] (0.2, 1.0) rectangle (11.6, 2.2);
    \draw[thick, dashed, targetgreen!80] (0.2, 2.2) -- (11.6, 2.2) node[right, font=\tiny\bfseries, text=targetgreen!90!black] {$L + \eps$};
    \draw[thick, sourceblue!90!black] (0.2, 1.6) -- (11.6, 1.6) node[right, font=\tiny\bfseries, text=sourceblue!90!black] {Limit $L$};
    \draw[thick, dashed, targetgreen!80] (0.2, 1.0) -- (11.6, 1.0) node[right, font=\tiny\bfseries, text=targetgreen!90!black] {$L - \eps$};
    
    % Threshold N line
    \draw[very thick, trapred!80, dashed] (5.0, 0.0) -- (5.0, 3.0) node[above, font=\scriptsize\bfseries, text=trapred!90!black] {Cutoff $N$};
    
    % Sequence points
    % Before N (oscillating outside band)
    \node[circle, fill=sourceblue, inner sep=1.8pt, label={above:\tiny $a_1$}] at (0.8, 2.8) {};
    \node[circle, fill=sourceblue, inner sep=1.8pt, label={below:\tiny $a_2$}] at (1.6, 0.4) {};
    \node[circle, fill=sourceblue, inner sep=1.8pt, label={above:\tiny $a_3$}] at (2.4, 2.5) {};
    \node[circle, fill=sourceblue, inner sep=1.8pt, label={below:\tiny $a_4$}] at (3.3, 0.7) {};
    \node[circle, fill=sourceblue, inner sep=1.8pt, label={above:\tiny $a_5$}] at (4.2, 2.3) {};
    
    % After N (trapped forever inside band)
    \node[circle, fill=targetgreen!80!black, inner sep=1.8pt, label={above:\tiny $a_N$}] at (5.3, 1.9) {};
    \node[circle, fill=targetgreen!80!black, inner sep=1.8pt] at (6.2, 1.35) {};
    \node[circle, fill=targetgreen!80!black, inner sep=1.8pt] at (7.1, 1.75) {};
    \node[circle, fill=targetgreen!80!black, inner sep=1.8pt] at (8.0, 1.48) {};
    \node[circle, fill=targetgreen!80!black, inner sep=1.8pt] at (8.9, 1.68) {};
    \node[circle, fill=targetgreen!80!black, inner sep=1.8pt] at (9.8, 1.55) {};
    \node[circle, fill=targetgreen!80!black, inner sep=1.8pt] at (10.7, 1.62) {};
    
    \node[font=\scriptsize, align=left, text=targetgreen!90!black] at (8.5, 0.3) {All terms trapped inside $\eps$-tube for $n \ge N$};
  \end{scope}

  % Lane 2: Function Epsilon-Delta Box
  \begin{scope}[yshift=-4.8cm]
    \node[anchor=west, font=\small\bfseries] at (-0.5, 3.8) {2. The Function Contract: $\forall \eps > 0, \exists \delta > 0, 0 < |x - c| < \delta \implies |f(x) - L| < \eps$};
    \draw[rounded corners=4pt, fill=gray!3, draw=gray!25] (-0.6, -0.6) rectangle (12.2, 3.6);
    
    % Axes
    \draw[->, thick, gray!60] (0.5, 0.2) -- (8.5, 0.2) node[right, font=\scriptsize] {$x$};
    \draw[->, thick, gray!60] (1.0, 0.0) -- (1.0, 3.2) node[above, font=\scriptsize] {$y$};
    
    % Shaded safe rectangle (delta x epsilon)
    \fill[targetgreen!15] (3.4, 0.8) rectangle (5.8, 2.4);
    \draw[thick, draw=targetgreen!70] (3.4, 0.8) rectangle (5.8, 2.4);
    
    % Limits and deltas
    % Horizontal epsilon band
    \draw[dashed, targetgreen!80] (1.0, 2.4) node[left, font=\tiny] {$L+\eps$} -- (8.0, 2.4);
    \draw[dashed, sourceblue!80] (1.0, 1.6) node[left, font=\tiny\bfseries] {$L$} -- (8.0, 1.6);
    \draw[dashed, targetgreen!80] (1.0, 0.8) node[left, font=\tiny] {$L-\eps$} -- (8.0, 0.8);
    
    % Vertical delta band
    \draw[dashed, quotientorange!80] (3.4, 0.2) node[below, font=\tiny] {$c-\delta$} -- (3.4, 2.8);
    \draw[dashed, sourceblue!80] (4.6, 0.2) node[below, font=\tiny\bfseries] {$c$} -- (4.6, 2.8);
    \draw[dashed, quotientorange!80] (5.8, 0.2) node[below, font=\tiny] {$c+\delta$} -- (5.8, 2.8);
    
    % Center point (c, L)
    \node[circle, fill=sourceblue, inner sep=2pt] at (4.6, 1.6) {};
    
    % Curve y = f(x) passing safely through vertical sides
    \draw[very thick, structurepurple] (2.0, 0.3) to[out=25, in=215] (3.4, 1.0) to[out=35, in=215] (4.6, 1.6) to[out=35, in=205] (5.8, 2.2) to[out=25, in=200] (7.5, 2.9);
    \node[structurepurple!90!black, font=\scriptsize\bfseries] at (7.8, 2.9) {$y=f(x)$};
    
    % Annotations
    \node[anchor=west, font=\scriptsize, align=left] at (8.7, 1.6) {
      $\bullet$ Skeptic dictates height $\pm\eps$.\\
      $\bullet$ Prover responds with width $\pm\delta$.\\
      $\bullet$ Curve exits through the \textbf{vertical}\\
      \quad boundaries, never the $\eps$-ceiling\\
      \quad or $\eps$-floor!
    };
  \end{scope}
\end{tikzpicture}
\end{center}

\subsection{Worked Proof: \texorpdfstring{$1/n \to 0$}{1/n -> 0}}

\begin{theorem}
\[
\lim_{n \to \infty} \frac{1}{n} = 0.
\]
\end{theorem}

\begin{proof}
Let $\eps > 0$ be arbitrary. 

\emph{Scratchpad Search:} We want $|a_n - L| < \eps$, which means:
\[
\left|\frac{1}{n} - 0\right| = \frac{1}{n} < \eps \iff n > \frac{1}{\eps}.
\]
This tells us exactly how far out in the sequence we need to travel.

\emph{Official Written Proof:}
By the Archimedean Property of real numbers, there exists a positive integer $N \in \N$ such that
\[
N > \frac{1}{\eps}.
\]
Now let $n \in \N$ be an arbitrary index and assume $n \ge N$.
Because $n \ge N$, taking reciprocals reverses the inequality:
\[
\frac{1}{n} \le \frac{1}{N}.
\]
Combining this with our choice of $N > 1/\eps$ (which means $1/N < \eps$):
\[
\left|\frac{1}{n} - 0\right| = \frac{1}{n} \le \frac{1}{N} < \eps.
\]
Thus for all $n \ge N$, $|1/n - 0| < \eps$. Since $\eps > 0$ was arbitrary, $\lim_{n \to \infty} 1/n = 0$.
\end{proof}

\subsection{The Continuous Contract: Epsilon and Delta}

The exact same game governs continuous functions, with one twist: instead of choosing an integer index cutoff $N$, you choose an input neighborhood radius $\delta > 0$.

\begin{definition}[Continuity at a Point]
A function $f: \R \to \R$ is \textbf{continuous at a point $c \in \R$} if:
\[
(\forall \eps > 0)(\exists \delta > 0)(\forall x \in \R)(|x - c| < \delta \implies |f(x) - f(c)| < \eps).
\]
\end{definition}

Geometrically, the skeptic draws a horizontal strip of half-height $\eps$ centered at $y = f(c)$. You respond by constructing a vertical window of half-width $\delta$ centered at $x = c$. You win if the graph of $f(x)$ stays trapped within the skeptic's strip across the entire width of your window.

\begin{theorem}
The linear function $f(x) = 3x + 2$ is continuous at every $c \in \R$.
\end{theorem}

\begin{proof}
Let $c \in \R$ and $\eps > 0$ be arbitrary.

\emph{Scratchpad Search:} Because $f(x) = 3x + 2$ has slope $3$, any horizontal input deviation $|x - c|$ is stretched vertically by a factor of $3$:
\[
|f(x) - f(c)| = |(3x + 2) - (3c + 2)| = 3|x - c| < \eps \iff |x - c| < \frac{\eps}{3}.
\]
This tells us to divide the output tolerance by the slope and set $\delta = \eps / 3$.

\emph{Official Proof:}
Choose $\delta = \frac{\eps}{3} > 0$. Suppose $|x - c| < \delta$. Then:
\[
|f(x) - f(c)| = |(3x + 2) - (3c + 2)| = 3|x - c| < 3\delta = 3\left(\frac{\eps}{3}\right) = \eps.
\]
Thus $f$ is continuous at $c$.
\end{proof}

\subsection{Why Analysis Is Proof Training}

Real analysis has a reputation for being demanding, but its difficulty rarely comes from complicated arithmetic. Its challenge lies in \textbf{quantifier order} and tracking which variables are allowed to depend on which earlier choices. Notice how for the linear function $f(x) = 3x + 2$, the slope is $3$ everywhere, so $\delta = \eps/3$ works at every point $c$. For a curved graph like $f(x) = x^2$, the slope $f'(c) = 2c$ grows steeper as $c$ moves to the right—meaning a fixed vertical band $\pm\eps$ forces a narrower and narrower horizontal window $\delta$ as $c$ increases, so $\delta$ must depend on both $\eps$ and $c$.

Notice the strict dependency chain:
\[
\forall \eps > 0 \quad \longrightarrow \quad \exists N \in \N \quad \longrightarrow \quad \forall n \ge N.
\]
Your choice of $N$ is permitted to depend on $\eps$. It is \emph{not} permitted to depend on $n$, because $n$ has not been introduced yet!

If you swap the order of the first two quantifiers:
\[
\exists N \in \N \quad \forall \eps > 0 \quad \forall n \ge N \dots
\]
you would be claiming that a single fixed cutoff $N$ works for \emph{every} tolerance simultaneously. If a sequence is trapped within $\eps = 10^{-100}$ past index $N$ for all $\eps$, every term past $N$ must equal $L$ exactly! Swapping two words in the quantifier prefix transforms a subtle statement about convergence into a trivial statement about sequences that become constant.

When you practice $\eps$-$N$ and $\eps$-$\delta$ proofs, you are not just learning analysis. You are training your logical reflexes to respect the scoping, dependencies, and rules of formal mathematical thought.

\begin{trapbox}
\textbf{Cauchy's Great Quantifier Slip (1821):}
Even the greatest pioneers stumbled when navigating quantifier order. In his 1821 textbook \emph{Cours d'Analyse}, Cauchy stated as an established theorem that if an infinite series of continuous functions converges, its sum must also be continuous. 
Five years later, Niels Henrik Abel and others showed this was false: a series of smooth trigonometric waves can converge to a square wave with vertical cliff edges and jump discontinuities! 
What had gone wrong? Cauchy had assumed that because each individual point $x$ eventually settles down ($\forall x \, \forall \eps > 0 \, \exists N$), the whole curve must settle down at a single shared speed ($\forall \eps > 0 \, \exists N \, \forall x$). He accidentally swapped the order of ``for every $x$'' and ``there exists a cutoff $N$''! It took over two decades for Philipp Ludwig von Seidel (1847) and George Gabriel Stokes to untangle this confusion, giving birth to the modern distinction between pointwise and uniform convergence.
\end{trapbox}

\subsection{What Analysis Unlocked, Its Limits, and Where We Go Next}

The $\eps$-$\delta$ revolution provided the rock-solid bedrock that saved calculus from collapse:
\begin{itemize}
\item \textbf{Calculus without contradiction:} Derivatives, integrals, and the Fundamental Theorem of Calculus were finally proved without invoking vanishing ghosts or dividing by zero.
\item \textbf{Infinite series discipline:} It allowed mathematicians to analyze power series, Fourier transforms, and wave equations without producing divergent absurdities.
\item \textbf{Logical precision:} It trained generations of mathematicians in quantifier scoping, variable binding, and adversarial proof construction.
\end{itemize}

Yet we must acknowledge where the $\eps$-$\delta$ contract encounters its boundary: it is fundamentally tethered to the **metric ruler of the real line**. 

The definitions of $|x - c| < \delta$ and $|f(x) - L| < \eps$ rely on measuring physical distances with absolute values. But in modern physics and geometry, we frequently study spaces where:
\begin{itemize}
\item There is no coordinate grid and no ruler measuring distances (such as curved spacetimes in general relativity).
\item Spaces can be bent, stretched, or compressed elastically without tearing (topology).
\item Spaces are defined algebraically by symmetries and group actions rather than numbers.
\end{itemize}
If you stretch a rubber sheet, every distance between points changes---meaning all $\eps$ and $\delta$ values are altered---yet the underlying continuity of the surface has not been broken at all! Tying continuity to a ruler is an artificial restriction.

This brings us to the threshold of advanced pure mathematics! In the chapters ahead:
\begin{itemize}
\item In Chapter~21, \textbf{Abstract Algebra Preview: Structure and Symmetry}, we step beyond numbers to study groups, rings, and fields---the algebraic languages of symmetry and lawful transformation.
\item In Chapter~22, \textbf{Topology Preview: Continuity Without Rulers}, we liberate continuity from $\eps$ and $\delta$ forever, learning how to define open sets and nearness using the pure preimage machinery of Chapter~14.
\item In Chapters~28 and 35, we will return to analysis for a deep second pass, exploring uniform convergence, Dedekind cuts, and compactness.
\end{itemize}

\subsection{Exercises}

\begin{enumerate}
\item Use the formal $\eps$-$N$ definition of a limit to prove:
\[
\lim_{n \to \infty} \frac{2n + 1}{n + 3} = 2.
\]
(Perform your scratchpad algebra first to isolate $|a_n - 2|$, bound the denominator cleanly, and then write the official proof.)
\item Prove that the sequence $a_n = (-1)^n$ does \emph{not} converge to $1$. 
(To negate the limit definition: show there exists an $\eps > 0$ such that for every $N \in \N$, there is some $n \ge N$ with $|a_n - 1| \ge \eps$.)
\item Use the $\eps$-$\delta$ definition to prove that $f(x) = x^2$ is continuous at $c = 3$. (Hint: Factor $|x^2 - 9| = |x - 3||x + 3|$, assume $\delta \le 1$ to bound $|x + 3|$, and choose $\delta = \min\{1, \eps/7\}$.)
\end{enumerate}

\newpage
